\documentclass[11pt]{article}
\usepackage[margin=1in]{geometry}
\usepackage{amsmath, amssymb}
\usepackage{xcolor}
\usepackage{tcolorbox}

\newcommand{\dd}[2]{\frac{d #1}{d #2}}

\title{The Exponential ODE: $\dd{u}{t} = \alpha\,u$\\[6pt]
\large Analytical Solution by Separation of Variables}
\author{Claude Opus 4.6}
\date{April 2026}

\begin{document}
\maketitle

%% ============================================================
\section{Problem Statement}

\begin{tcolorbox}[colback=blue!5, colframe=blue!50!black, title=The Exponential ODE]
Find $u(t)$ satisfying:
\begin{align}
  \textbf{ODE:}& \quad \dd{u}{t} = \alpha\,u, \qquad \alpha \in \mathbb{R} \label{eq:ode}\\[6pt]
  \textbf{IC:}& \quad u(0) = u_0 \label{eq:ic}
\end{align}
where $\alpha$ is a constant parameter and $u_0$ is the initial condition.
\end{tcolorbox}

\medskip
\noindent
This is the simplest first-order linear ODE with constant coefficients.
Despite its simplicity, it is the prototype for exponential growth ($\alpha > 0$),
exponential decay ($\alpha < 0$), and equilibrium ($\alpha = 0$).

%% ============================================================
\section{Solution by Separation of Variables}

\subsection{Step 1: Separate the variables}

Starting from \eqref{eq:ode}, move all $u$-dependence to the left and all $t$-dependence to the right.
Assuming $u \neq 0$, divide both sides by $u$:
\begin{equation}
  \frac{1}{u}\,\dd{u}{t} = \alpha
  \label{eq:separated_implicit}
\end{equation}

Now treat $\dd{u}{t}$ as a ratio of differentials and separate:
\begin{equation}
  \boxed{\frac{du}{u} = \alpha\,dt}
  \label{eq:separated}
\end{equation}

\subsection{Step 2: Integrate both sides}

Integrate the left side with respect to $u$ and the right side with respect to $t$:
\begin{align}
  \int \frac{du}{u} &= \int \alpha\,dt \\[6pt]
  \ln|u| &= \alpha\,t + C
  \label{eq:log_form}
\end{align}
where $C$ is the constant of integration.

\subsection{Step 3: Exponentiate to solve for $u$}

Exponentiate both sides of \eqref{eq:log_form}:
\begin{align}
  |u| &= e^{\alpha t + C} \\[6pt]
      &= e^C \cdot e^{\alpha t}
\end{align}

Since $e^C > 0$ is an arbitrary positive constant, and absorbing the sign from the absolute value,
we write $A = \pm e^C$ as a single arbitrary constant:
\begin{equation}
  u(t) = A\,e^{\alpha t}
  \label{eq:general}
\end{equation}

This is the \textbf{general solution} --- it contains one arbitrary constant $A$,
as expected for a first-order ODE.

\subsection{Step 4: Apply the initial condition}

Evaluate \eqref{eq:general} at $t = 0$ and use \eqref{eq:ic}:
\begin{align}
  u(0) &= A\,e^{\alpha \cdot 0} \\[6pt]
       &= A\,e^0 \\[6pt]
       &= A \cdot 1 \\[6pt]
       &= A
\end{align}

Therefore $A = u_0$, and the particular solution is:

\begin{tcolorbox}[colback=green!5, colframe=green!50!black, title=Analytical Solution (General IC)]
\begin{equation}
  \boxed{u(t) = u_0\,e^{\alpha t}}
  \label{eq:solution}
\end{equation}
\end{tcolorbox}

%% ============================================================
\section{Verification}

A solution is only valid if it satisfies both the ODE and the initial condition.

\subsection{Check 1: Does it satisfy the ODE?}

Differentiate \eqref{eq:solution} with respect to $t$:
\begin{equation}
  \dd{u}{t} = u_0 \cdot \alpha\,e^{\alpha t} = \alpha\,\bigl(u_0\,e^{\alpha t}\bigr) = \alpha\,u \quad\checkmark
\end{equation}

\subsection{Check 2: Does it satisfy the IC?}

Evaluate at $t = 0$:
\begin{equation}
  u(0) = u_0\,e^{\alpha \cdot 0} = u_0 \cdot 1 = u_0 \quad\checkmark
\end{equation}

Both conditions are satisfied. The solution is verified.

%% ============================================================
\section{Specific Case: $u_0 = \tfrac{1}{2}$, $t \in [0,\,1]$}

\begin{tcolorbox}[colback=red!5, colframe=red!50!black, title=Specific Problem]
\begin{align}
  \textbf{ODE:}& \quad \dd{u}{t} = \alpha\,u \\[6pt]
  \textbf{IC:}& \quad u(0) = \tfrac{1}{2} \\[6pt]
  \textbf{Domain:}& \quad t \in [0,\,1]
\end{align}
\end{tcolorbox}

\medskip
\noindent
Substituting $u_0 = \frac{1}{2}$ into \eqref{eq:solution}:

\begin{equation}
  \boxed{u(t) = \frac{1}{2}\,e^{\alpha t}, \qquad t \in [0,\,1]}
  \label{eq:specific}
\end{equation}

\medskip
\noindent
\textbf{Boundary values:}
\begin{align}
  u(0) &= \tfrac{1}{2}\,e^0 = \tfrac{1}{2} \\[6pt]
  u(1) &= \tfrac{1}{2}\,e^{\alpha}
\end{align}

\medskip
\noindent
\textbf{Behavior depends on the sign of $\alpha$:}

\begin{center}
\renewcommand{\arraystretch}{1.5}
\begin{tabular}{l|c|c|l}
  \hline
  \textbf{Regime} & $\alpha$ & $u(1)$ & \textbf{Behavior} \\
  \hline
  Growth    & $\alpha > 0$ & $> \frac{1}{2}$ & $u$ increases exponentially from $\frac{1}{2}$ \\
  Equilibrium & $\alpha = 0$ & $= \frac{1}{2}$ & $u$ remains constant at $\frac{1}{2}$ \\
  Decay     & $\alpha < 0$ & $< \frac{1}{2}$ & $u$ decreases exponentially toward $0$ \\
  \hline
\end{tabular}
\end{center}

\medskip
\noindent
Some representative values at $t = 1$:

\begin{center}
\renewcommand{\arraystretch}{1.4}
\begin{tabular}{c|c|c}
  \hline
  $\alpha$ & $u(1) = \frac{1}{2}e^{\alpha}$ & Decimal \\
  \hline
  $-2$  & $\frac{1}{2}e^{-2}$  & $\approx 0.0677$ \\
  $-1$  & $\frac{1}{2}e^{-1}$  & $\approx 0.1839$ \\
  $0$   & $\frac{1}{2}$        & $= 0.5000$ \\
  $1$   & $\frac{1}{2}e$       & $\approx 1.3591$ \\
  $2$   & $\frac{1}{2}e^{2}$   & $\approx 3.6945$ \\
  \hline
\end{tabular}
\end{center}

%% ============================================================
\section{Key Observations}

\begin{enumerate}
  \item \textbf{The exponential is its own derivative.}
    The function $e^{\alpha t}$ is the unique (up to scaling) eigenfunction of the
    differentiation operator $\frac{d}{dt}$ with eigenvalue $\alpha$.
    This is \emph{why} the exponential solves this ODE.

  \item \textbf{Separation of variables works because the ODE is separable.}
    The right-hand side $\alpha u$ factors as a function of $u$ alone times a
    function of $t$ alone (with the $t$-part being the constant $\alpha$).
    Not all ODEs have this structure.

  \item \textbf{This ODE is the building block for linear systems.}
    For a system $\frac{d\mathbf{u}}{dt} = A\,\mathbf{u}$ with constant matrix $A$,
    the solution is $\mathbf{u}(t) = e^{At}\,\mathbf{u}_0$. Each eigenvector of $A$
    evolves independently as $e^{\lambda_i t}$ --- i.e., each mode satisfies exactly
    this scalar ODE with $\alpha = \lambda_i$.

  \item \textbf{Stability is determined entirely by $\alpha$.}
    If $\alpha < 0$: the origin is a stable equilibrium (perturbations decay).
    If $\alpha > 0$: the origin is unstable (perturbations grow).
    If $\alpha = 0$: neutral stability (perturbations persist but do not grow).

  \item \textbf{The $u = 0$ edge case.}
    In Step~1 we divided by $u$, requiring $u \neq 0$.
    But $u(t) = 0$ for all $t$ is itself a valid solution (the trivial solution).
    Since our IC has $u_0 = \frac{1}{2} \neq 0$, the exponential solution is the
    correct one, and $u(t) > 0$ for all $t$.
\end{enumerate}

\end{document}
